%%%%%%%%%%%%%%%%%%%%%%%%%%%%%%%%%%%%%%%%%%%%%%%%%%%%%%%%%%%%%%%%%%%%%%%%%%%%
% Supporting Information for:
% "Pore-Scale Phase Trapping as a Mechanism for a Self-Locking Dissociation
%  Front in Hydrate-Bearing Sediments"
% Water Resources Research
%
% AGU Supporting Information template -- uses agutexSI2019.cls
% IMPORTANT: This file is NOT copyedited before publication.
% All references must also appear in the main manuscript reference list.
%%%%%%%%%%%%%%%%%%%%%%%%%%%%%%%%%%%%%%%%%%%%%%%%%%%%%%%%%%%%%%%%%%%%%%%%%%%%

\documentclass[draft,wrr]{agutexSI2019}

\usepackage{graphicx}
\setkeys{Gin}{draft=false}   % show figures even in draft mode
\graphicspath{{figures_si/}}

\usepackage{amsmath}
\usepackage{amssymb}
\usepackage{bm}
\usepackage{upgreek}  % provides \upmu for micrometer symbol

\usepackage{booktabs}
\usepackage{threeparttable}
\usepackage{multirow}
\usepackage{tabularx}
\usepackage{makecell}
\usepackage{array}
\usepackage{lineno}      % provides linenomath* environment for equations in draft mode
\usepackage{xcolor}
% Keep draft double-spacing but suppress draft timestamp in page footer.
\makeatletter
\def\@oddfoot{\jfootline}
\def\@evenfoot{\jfootline}
\makeatother

% Custom column types for tabularx
\newcolumntype{P}[1]{>{\raggedright\arraybackslash}p{#1}}
\newcolumntype{Y}{>{\raggedright\arraybackslash}X}

% Running heads (optional)
% \authorrunninghead{CHEN AND LI}
% \titlerunninghead{SUPPORTING INFORMATION}

% Corresponding author
% \authoraddr{Corresponding author: Shengli Li,
% College of Construction Engineering, Jilin University,
% Changchun, 130026, China.(lishengli@jlu.edu.cn)}

\begin{document}
\raggedbottom\relax
\pagestyle{headings}

%% -----------------------------------------------------------------------
%%  TITLE
%% -----------------------------------------------------------------------
\title{Supporting Information for ``Pore-Scale Phase Trapping in Hydrate-Bearing
Sediments: Implications for a Self-Locking Dissociation Front''}

%% -----------------------------------------------------------------------
%%  AUTHORS AND AFFILIATIONS
%% -----------------------------------------------------------------------
\authors{Yuxiang Chen\affil{1,2},
         Shengli Li\affil{1,2},
         Haonan Wu\affil{1,2},
         Qiang Liu\affil{1,2},
         Shijing Zhang\affil{1,2}}

\affiliation{1}{State Key Laboratory of Deep Earth Exploration and Imaging,
College of Construction Engineering, Jilin University, Changchun, 130026, China}
\affiliation{2}{Key Laboratory of Drilling and Exploitation Technology in Complex
Conditions of Natural Resources, College of Construction Engineering,
Jilin University, Changchun, 130026, China}

%% -----------------------------------------------------------------------
%%  BEGIN ARTICLE
%% -----------------------------------------------------------------------
\begin{article}

%% -----------------------------------------------------------------------
%%  CONTENTS OF THIS FILE
%% -----------------------------------------------------------------------
\noindent\textbf{Contents of this file}
\begin{enumerate}
  \item Text S1: Multi-component pseudopotential lattice Boltzmann method
  \item Text S2: Model validation
  \item Text S3: Validation of the trapping criterion
  \item Text S4: Model Validation and SHAP Attribution Analysis
  \item Figures S1 to S10
  \item Tables S1 to S2
\end{enumerate}


%% -----------------------------------------------------------------------
%%  INTRODUCTION
%% -----------------------------------------------------------------------
\noindent\textbf{Introduction}

This Supporting Information provides the full formulations, parameter
definitions, and implementation details of the multicomponent multi-relaxation-time
(MRT) pseudopotential lattice Boltzmann (LBM) framework used in the main
manuscript. It also documents the physical scaling procedure, the trapping
characterization protocol, the simulation design, and the model validation
benchmarks (Laplace pressure, contact angle, two-phase Poiseuille flow, and mesh
independence). Section S3 further details the validation of the trapping criterion
used to classify trapped clusters, and Section S4 describes the model validation
protocol and the SHAP attribution analysis framework.

%% ----------------------------------------------------------------------
%%  TEXT S1 -- MRT-LBM Framework
%% -----------------------------------------------------------------------
\noindent\textbf{Text S1.}
\textbf{Multicomponent Pseudopotential Lattice Boltzmann Framework}

\noindent\textit{S1.1 Governing Evolution Equation}

This study employs a multicomponent pseudopotential lattice Boltzmann framework
based on the multi-relaxation-time (MRT) formulation to simulate capillary-driven
gas-water flow in hydrate-bearing porous media. Compared with the conventional
single-relaxation-time (BGK) scheme, the MRT approach improves numerical stability
and accuracy under large viscosity and density contrasts, which is critical for
reproducing realistic interfacial dynamics and capillary phenomena in heterogeneous
domains.

The evolution of the particle distribution function for each fluid component
follows the MRT lattice Boltzmann equation:
\begin{linenomath*}
\begin{equation}
\label{eq:mrt-evolution}
f_{k,\alpha}(\mathbf{x}+\mathbf{e}_{\alpha}\delta_t,\, t+\delta_t)
=
f_{k,\alpha}(\mathbf{x},t)
-
\bar{\Lambda}_{\alpha\beta}
\left(
f_{k,\beta}-f_{k,\beta}^{\mathrm{eq}}
\right)\big|_{(\mathbf{x},t)}
+
\delta_t\, F'_{k,\alpha}\big|_{(\mathbf{x},t)} \, .
\end{equation}
\end{linenomath*}
where $f_{\alpha}$ is the density distribution function (DF),
$f_{\alpha}^{\mathrm{eq}}$ represents the equilibrium DF,
$F'_{\alpha}$ denotes the external pseudopotential force,
$\mathbf{e}_{\alpha}$ are the discrete velocities:

\begin{linenomath*}
\begin{equation}
\label{eq:discrete-velocities}
\mathbf{e}_{\alpha} = \begin{cases}
(0,0) & \alpha=0 \\[1.5ex] % 这里增加了 1.5ex 的额外垂直间距
c\left(\cos \left[\frac{(\alpha-1) \pi}{2}\right], \sin \left[\frac{(\alpha-1) \pi}{2}\right]\right) & \alpha=1 \sim 4 \\[1.5ex]
\sqrt{2} c\left(\cos \left[\frac{(2 \alpha-9) \pi}{4}\right], \sin \left[\frac{(2 \alpha-9) \pi}{4}\right]\right) & \alpha=5 \sim 8
\end{cases}
\end{equation}
\end{linenomath*}
The spatial and temporal steps $\delta_x$ and $\delta_t$
are both set to unity. The two-dimensional nine-velocity (D2Q9) lattice model
is adopted, with lattice constant $c = \delta_x/\delta_t$ and sound speed
$c_s = c/\sqrt{3}$.
$\Lambda_{\alpha\beta} = {M}^{-1}{\Lambda}{M}$ is the collision matrix
constructed from the moment-space transform $M$ and the diagonal relaxation
matrix $\Lambda$. The matrices $M$ and $\Lambda$ are defined as~\cite{lallemand_theory_2000}:

\begin{linenomath*}
\begin{equation}
\label{eq:transform-matrix}
\mathbf{M} = \begin{bmatrix}
1 & 1 & 1 & 1 & 1 & 1 & 1 & 1 & 1 \\
-4 & -1 & -1 & -1 & -1 & 2 & 2 & 2 & 2 \\
4 & -2 & -2 & -2 & -2 & 1 & 1 & 1 & 1 \\
0 & 1 & 0 & -1 & 0 & 1 & -1 & -1 & 1 \\
0 & -2 & 0 & 2 & 0 & 1 & -1 & -1 & 1 \\
0 & 0 & 1 & 0 & -1 & 1 & 1 & -1 & -1 \\
0 & 0 & -2 & 0 & 2 & 1 & 1 & -1 & -1 \\
0 & 1 & -1 & 1 & -1 & 0 & 0 & 0 & 0 \\
0 & 0 & 0 & 0 & 0 & 1 & -1 & 1 & -1
\end{bmatrix}
\end{equation}
\end{linenomath*}

\begin{linenomath*}
\begin{equation}
\label{eq:relaxation-matrix}
\Lambda = \mathrm{diag} \left( \tau_{\rho}^{-1}, \tau_{e}^{-1}, \tau_{\zeta}^{-1}, \tau_{j}^{-1}, \tau_{q}^{-1}, \tau_{j}^{-1}, \tau_{q}^{-1}, \tau_{v}^{-1}, \tau_{v}^{-1} \right)
\end{equation}
\end{linenomath*}
$\tau_{v}$ and $\tau_{e}$ determine the kinematic viscosity and bulk viscosity, respectively,
with $\nu = c_{s}^{2}(\tau_{v}-0.5)$ and $\zeta = c_{s}^{2}(\tau_{e}-0.5)$.
The specific values related to $\Lambda$ are specified as follows~\cite{hou_multi_2018}:
$\tau_{\rho}^{-1} = \tau_{j}^{-1} = 1$;
$\tau_{e}^{-1} = \tau_{\zeta}^{-1} = 0.8$; and $\tau_{q}^{-1} = 1.1$.



Using the transformation matrix $M$, both $f$ and $f^{\mathrm{eq}}$ are
projected into the moment space through $m = Mf$ and $m^{\mathrm{eq}} = Mf^{\mathrm{eq}}$.
The right-hand side of equation~\eqref{eq:mrt-evolution} can then be rewritten as:
\begin{linenomath*}
\begin{equation}
\label{eq:mrt-moment-update}
\mathbf{m}^{*}=\mathbf{m}-\bm{\Lambda}\left(\mathbf{m} - \mathbf{m}^{\mathrm{eq}}\right)
+ \delta_t\left(\mathbf{I} - \frac{\bm{\Lambda}}{2}\right)\bar{\mathbf{S}}+\delta_t\,\mathbf{C} \, .
\end{equation}
\end{linenomath*}
In equation~\eqref{eq:mrt-moment-update}, the collision operator acts in moment
space, where the diagonal relaxation matrix $\Lambda$ assigns independent
relaxation rates to individual kinetic moments. This freedom allows non-conserved
higher-order moments to relax more rapidly, suppressing spurious velocities and
improving numerical stability under strong density and viscosity contrasts.

The macroscopic density $\rho_k$ and velocity $\mathbf{u}$ are defined as:
\begin{linenomath*}
\begin{equation}
\rho_k = \sum_{\alpha} f_{k,\alpha},
\qquad
\mathbf{u}
=
\frac{
\sum_{k}
\left(
\sum_{\alpha}
\mathbf{e}_{\alpha} f_{k,\alpha}
+
\frac{\delta_t}{2}\mathbf{F}_k
\right)
}{
\sum_{k} \rho_k
}
\end{equation}
\end{linenomath*}
where $\mathbf{F}_{k} = \mathbf{F}_{k,m} + \mathbf{F}_{k,\text{ads}} + \mathbf{F}_{b}$
is the total force on component $k$, including fluid--fluid intermolecular force,
fluid--solid interaction force, and the imposed body force.

A uniform body force $\mathbf{F}_b = \rho_k \mathbf{g}$ is applied along the
flow direction to drive gas--water displacement~\cite{ji_relative_2022}, where $g$ denotes the forcing
intensity. This forcing mimics pressure-gradient conditions under depressurization
and provides a controlled means of probing capillary-viscous competition.

\noindent\textit{S1.2 Interfacial Interactions and Phase Separation}

The multicomponent pseudopotential formulation governs the capillary forces that
drive phase segregation between water and methane~\cite{chen_critical_2014}.
In this model, the fluid-fluid interaction acting on component $k$ combines an
intra-component term $\mathbf{F}_k^{\mathrm{intra}}$ and an inter-component term
$\mathbf{F}_k^{\mathrm{inter}}$:
\begin{linenomath*}
\begin{equation}
\label{eq:force_total}
\mathbf{F}_{k,m}(\mathbf{x})
=
\mathbf{F}_k^{\mathrm{intra}}(\mathbf{x})
+
\mathbf{F}_k^{\mathrm{inter}}(\mathbf{x}) .
\end{equation}
\end{linenomath*}

\begin{linenomath*}
\begin{equation}
\label{eq:force_intra}
\mathbf{F}_k^{\mathrm{intra}}(\mathbf{x})
=
-
G_{kk}
\,
\psi_k(\mathbf{x})
\sum_{\alpha}
w(|\mathbf{e}_\alpha|^2)
\,
\psi_k(\mathbf{x}+\mathbf{e}_\alpha)
\,
\mathbf{e}_\alpha .
\end{equation}
\end{linenomath*}

\begin{linenomath*}
\begin{equation}
\label{eq:force_inter}
\mathbf{F}_k^{\mathrm{inter}}(\mathbf{x})
=
-
G_{k\bar{k}}
\,
\phi_k(\mathbf{x})
\sum_{\alpha}
w_\alpha
\,
\phi_{\bar{k}}(\mathbf{x}+\mathbf{e}_\alpha)
\,
\mathbf{e}_\alpha .
\end{equation}
\end{linenomath*}
Here $G_{kk}$ and $G_{k\bar{k}}$ denote the interaction strengths between the
same and different components, respectively. In this work, water (component 1)
and methane (component 2) interact through a tunable coupling strength $G_{12}$,
which controls gas--water phase separation. The weight coefficients $w_\alpha$
equal $4/9$ (rest), $1/9$ (face), and $1/36$ (diagonal) in the D2Q9 model.
Different from the pseudopotential  $\psi_k(\mathbf{x})$ in equation~\eqref{eq:force_intra},
an empirical density-based pseudopotential $\phi_k(\mathbf{x})$ is employed
in equation~\eqref{eq:force_inter} to ensure a proper component distribution~\cite{zheng_single_2019}.

To represent the strong equation-of-state nonlinearity of water, a
PR-EOS-based pseudopotential $\psi_1(\mathbf{x})$ is employed~\cite{he_thermodynamic_xxxx,yuan_equations_2006}:
\begin{linenomath*}
\begin{equation}
\label{eq:psi_water}
\psi_1(\mathbf{x})
=
\sqrt{
\frac{
2\left(
p_{k,\mathrm{EOS}}
-
\rho_1 c_s^2
\right)
}{
c^2 G_{11}
}
} .
\end{equation}
\end{linenomath*}
The Peng--Robinson equation~\cite{yuan_equations_2006} of state determines the local pressure $p_{k,\mathrm{EOS}}$:
\begin{linenomath*}
\begin{equation}
\label{eq:pr_eos}
p_{k,\mathrm{EOS}}
=
\frac{\rho_1 R T}{1 - b\rho_1}
-
\frac{a\,\alpha(T)\,\rho_1^2}{1 + 2b\rho_1 - b^2\rho_1^2}
\end{equation}
\end{linenomath*}

\begin{linenomath*}
\begin{equation}
\label{eq:phi-T}
\alpha(T) = \left[ 1 + (0.37464 + 1.54226\omega - 0.26992\omega^2) (1 - \sqrt{T/T_c}) \right]^2
\end{equation}
\end{linenomath*}
where $\omega = 0.344$, $a = 0.45724 \frac{R^2 T_c^2}{P_c}$, $b = 0.0778 \frac{R T_c}{P_c}$, $T_c$, $P_c$ are the
critical temperature and pressure, respectively.
In this work, the parameters are specified as $a = 1/49$, $b = 2/21$, and $R = 1$.

Methane is treated using an empirical density-based pseudopotential~\cite{xin_effect_2021}:
\begin{linenomath*}
\begin{equation}
\label{eq:psi_methane}
\psi_2(\mathbf{x})
=
\rho_0
\left[
1
-
\exp
\left(
-
\frac{\rho_2(\mathbf{x})}{\rho_0}
\right)
\right]
\end{equation}
\end{linenomath*}
where $\rho_0 = 1$ is a constant coefficient. This asymmetric treatment allows
water to retain a realistic capillary response while keeping methane numerically
tractable under large density ratios.

\noindent\textit{S1.3 Improved Forcing Scheme}

Under large density contrasts, standard forcing schemes introduce spurious
pressure shifts. To maintain thermodynamic consistency and suppress unphysical
currents along curved interfaces, we adopt the improved forcing scheme proposed
by \citeA{li_lattice_2013}, in which the force term is projected into moment
space and scaled by a parameter $\sigma$:
\begin{linenomath*}
\begin{equation}
\label{eq:moment_force}
\bar{\mathbf{S}}
=
\begin{bmatrix}
0 \\
6\left(u_x F_{k,x} + u_y F_{k,y}\right)
+
\displaystyle
\frac{
12\sigma
\left|\mathbf{F}_{k,m}\right|^2
}{
\psi^2 \Delta t \left(\tau_e - 0.5\right)
}
\\[18pt]
-6\left(u_x F_{k,x} + u_y F_{k,y}\right)
-
\displaystyle
\frac{
12\sigma
\left|\mathbf{F}_{k,m}\right|^2
}{
\psi^2 \Delta t \left(\tau_s - 0.5\right)
}
\\[18pt]
F_{k,x} \\
- F_{k,x} \\
F_{k,y} \\
- F_{k,y} \\
2\left(u_x F_{k,x} - u_y F_{k,y}\right) \\
u_x F_{k,y} + u_y F_{k,x}
\end{bmatrix} .
\end{equation}
\end{linenomath*}
The parameter $\sigma$ is tuned to achieve thermodynamic consistency; for water,
$\sigma = 0.11$ follows \citeA{hou_multi_2018}, while for methane $\sigma = 0$.

\noindent\textit{S1.4 Wettability Representation}

Wall wettability is represented by a fluid-solid interaction force
$\mathbf{F}_{k,\mathrm{ads}}$, which prescribes the contact angle between
fluids and solid surfaces. A modified pseudopotential approach~\cite{yang_improved_2021}
is adopted to impose a broad range of contact angles while maintaining stable
interfacial representation:
\begin{linenomath*}
\begin{equation}
\label{eq:force_wettability}
\mathbf{F}_{k,\mathrm{ads}}(\mathbf{x})
=
-
G_{k,\mathrm{ads}}
\,
\psi_k(\mathbf{x})
\sum_i
w\!\left(|\mathbf{e}_i|^2\right)
\,
\psi(\mathbf{x})
\,
s(\mathbf{x}+\mathbf{e}_i \Delta t)
\,
\mathbf{e}_i .
\end{equation}
\end{linenomath*}
Here $G_{k,\mathrm{ads}}$ is the adhesive coefficient used to tune wettability,
and the switch function $s(\mathbf{x})$ equals 1 at solid nodes and 0 at fluid
nodes. A ghost layer near the solid boundary is employed, with its density
estimated from neighboring fluid nodes~\cite{yang_improved_2021,li_implementation_2019}.
 \begin{linenomath*}
\begin{equation}
\label{eq:average-density}
\rho_{\mathrm{ave}}(\mathbf{x}) = \frac{\displaystyle \sum_{\alpha}
w(|\mathbf{e}_{\alpha}|^2) \rho(\mathbf{x} + \mathbf{e}_{\alpha} \delta_t)
[1 - s(\mathbf{x} + \mathbf{e}_{\alpha})]} {\displaystyle \sum_{\alpha}
w(|\mathbf{e}_{\alpha}|^2) [1 - s(\mathbf{x} + \mathbf{e}_{\alpha})]}
\end{equation}
\end{linenomath*}
The adhesive coefficient $G_{k,\mathrm{ads}}$
maps directly and linearly to the static contact angle $\theta$, as confirmed by
the contact angle benchmark (Figure~S2).


\noindent\textit{S1.5 Surface Tension Control}

In the original pseudopotential model, surface tension is intrinsically coupled
with other fluid properties and cannot be tuned independently. To introduce
explicit control over capillarity, an additional source term $\mathbf{C}$ is
incorporated following \citeA{li_achieving_2013}, allowing the interfacial force
to be adjusted without affecting bulk thermodynamic behavior:
\begin{linenomath*}
\begin{equation}
\label{eq:surface_tension_source}
\mathbf{C}
=
\left[
\renewcommand{\arraystretch}{1.6}
\begin{array}{c}
\makebox[3cm]{0} \\
1.5\,\tau_e^{-1}\left(Q_{xx}+Q_{yy}\right) \\
-1.5\,\tau_s^{-1}\left(Q_{xx}+Q_{yy}\right) \\
0 \\
0 \\
0 \\
0 \\
-\tau_v^{-1}\left(Q_{xx}-Q_{yy}\right) \\
-\tau_v^{-1} Q_{xy}
\end{array}
\right] .
\end{equation}
\end{linenomath*}
The variables $Q_{xx}$, $Q_{yy}$, and $Q_{xy}$ are obtained from the tensor
$\mathbf{Q}$:
\begin{linenomath*}
\begin{equation}
\label{eq:Q_tensor}
\mathbf{Q}
=
\kappa
\,
\frac{G_{11}}{2}
\,
\psi_1(\mathbf{x})
\sum_{\alpha}
w\!\left(|\mathbf{e}_{\alpha}|^2\right)
\left[
\psi_1(\mathbf{x}+\mathbf{e}_{\alpha})
-
\psi_1(\mathbf{x})
\right]
\,
\mathbf{e}_{\alpha}\mathbf{e}_{\alpha} .
\end{equation}
\end{linenomath*}
The parameter $\kappa \in [0,1]$ enables independent adjustment of the surface
tension over a broad range.
The system pressure for the multi-component LB model can be calculated by:
\begin{linenomath*}
\begin{equation}
\label{eq:system-pressure}
p = c_{s}^{2} \sum_{k} \rho_{k} + 0.5 c^{2} \sum_{k} G_{kk}
\psi_{k}^{2} + 0.5 c^{2} \sum_{k \neq \bar{k}} G_{k\bar{k}} \phi_k \phi_{\bar{k}}
\end{equation}
\end{linenomath*}
%% -----------------------------------------------------------------------
%%  TEXT S2 -- Physical Scaling and Simulation Parameters
%% -----------------------------------------------------------------------
\noindent\textbf{Text S2.}
\textbf{Model Validation Benchmarks}

The MRT-based pseudopotential model was validated through four standard
benchmarks to confirm its ability to reproduce interfacial tension, wettability
control, phase-mobility partitioning, and spatial resolution sufficiency.

\noindent\textit{S2.1 Laplace Pressure Test}

A droplet with various initial radii $R$ was placed at the center of a
$300 \times 300$ lattice domain with periodic boundaries. The simulation temperature
was set to $0.85\,T_c$ and the surface tension was tuned via $\kappa$ to match
target values for each equilibrium state. Figure~S1 shows the linear dependence of
the pressure difference $\Delta P$ on $1/R$ for both equilibrium states, consistent
with the Laplace relation $\Delta P = \sigma / R$. The strong linearity
($R^2 \approx 0.99$) confirms that the model accurately reproduces the Laplace law.

\noindent\textit{S2.2 Contact Angle Test}

The wettability boundary condition was validated by simulating a droplet resting on
a flat solid surface within a $300 \times 100$ lattice domain. Bounce-back and
wettability boundary conditions were applied at the bottom wall and periodic
boundaries elsewhere. A droplet of radius 30 lattice units was allowed to relax to
its equilibrium configuration. Figure~S2 illustrates the equilibrium droplets under
both equilibrium states for various $G_{k,\mathrm{ads}}$ values. A clear linear
relationship between the resulting contact angle $\theta$ and $G_{k,\mathrm{ads}}$
confirms that the adhesive coefficient can be directly mapped to the static contact
angle.

\noindent\textit{S2.3 Two-Dimensional Two-Phase Poiseuille Flow}

To further validate the two-phase flow model, a two-dimensional Poiseuille flow of
immiscible fluids in a channel was simulated and compared with the analytical
solution derived from the Navier--Stokes equations. The flow configuration and
phase distribution are illustrated in Figure~S3, which schematically shows the
wetting and non-wetting regions and the coordinate system adopted in the analysis.
The analytical solution for the velocity profile is expressed as:
\begin{linenomath*}
\begin{equation}
u_x(y)=
\begin{cases}
A_{\mathrm{nw}}y^2 + C_{\mathrm{nw}}, & |y|\le a \\
A_{\mathrm{w}}y^2 + B_{\mathrm{w}}y + C_{\mathrm{w}}, & a \le |y|\le b
\end{cases}
\label{eq:poiseuille_velocity_profile}
\end{equation}
\end{linenomath*}
where
\begin{linenomath*}
\begin{equation}
A_{\mathrm{nw}}=\frac{F_b}{\mu_{\mathrm{nw}}},\quad
A_{\mathrm{w}}=\frac{F_b}{\mu_{\mathrm{w}}},\quad
B_{\mathrm{w}}=-2A_{\mathrm{w}}a+2MA_{\mathrm{nw}}a,
\label{eq:poiseuille_coeff_1}
\end{equation}
\end{linenomath*}
\begin{linenomath*}
\begin{equation}
C_{\mathrm{nw}}=(A_{\mathrm{w}}-A_{\mathrm{nw}})a^2-B_{\mathrm{w}}(b-a)-A_{\mathrm{w}}b^2,\quad
C_{\mathrm{w}}=-A_{\mathrm{w}}b^2-B_{\mathrm{w}}b.
\label{eq:poiseuille_coeff_2}
\end{equation}
\end{linenomath*}
Here $F_b$ is the uniform body force applied along the $x$-direction, $a$ and
$b$ denote the half-widths of the non-wetting and wetting regions,
$\mu_{\mathrm{nw}}$ and $\mu_{\mathrm{w}}$ are the dynamic viscosities of the
non-wetting and wetting fluids, respectively, and
$M=\mu_{\mathrm{nw}}/\mu_{\mathrm{w}}$ is the viscosity ratio.

From the analytical velocity profile, $M$ and the phase saturations jointly
determine the flow resistance and permeability partition between phases. The
analytical relative permeabilities for wetting and non-wetting phases~\cite{liang_analytical_2018} are:
\begin{linenomath*}
\begin{equation}
k_{r,w}=\frac{1}{2}S_w^2(3-S_w)
\label{eq:poiseuille_krw}
\end{equation}
\end{linenomath*}
\begin{linenomath*}
\begin{equation}
k_{r,\mathrm{nw}}=S_{\mathrm{nw}}\left[\frac{3}{2}M+S_{\mathrm{nw}}^2\left(1-\frac{3}{2}M\right)\right]
\label{eq:poiseuille_krnw}
\end{equation}
\end{linenomath*}
where $S_w$ and $S_{\mathrm{nw}}=1-S_w$ are the saturations of the wetting and
non-wetting phases, respectively.

To reproduce this flow configuration numerically, a $300\times200$ lattice
domain was employed with no-slip boundaries at the top and bottom and periodic
boundaries at the lateral sides. The co-current two-phase flow was driven by a
uniform body force ($F_b = 5 \times 10^{-9}$). Figure~S4 presents the validation
results for two equilibrium states (Eq1 and Eq2), showing the comparison between
LBM simulations and analytical solutions for both gas and water phases. The
simulated relative permeabilities show excellent agreement with the analytical
predictions ($R^2 \approx 0.99$), confirming the robustness and reliability of the
proposed model.

\noindent\textit{S2.4 Mesh Independence Study}

To verify the reliability of the numerical model and determine an appropriate
spatial resolution, a mesh-independence assessment was conducted based on the
Poiseuille-flow benchmark described in Section~S2.3. Four nested uniform grids
with $N_Y=100$, $200$, $400$, and $800$ were evaluated. During the tests, the
external driving force was adjusted to maintain a constant flow rate and thus a
consistent Reynolds number across all grid configurations.

The relative error between the numerically predicted flow rate and the
analytical solution is defined as:
\begin{linenomath*}
\begin{equation}
\varepsilon=
\left|
\frac{Q_{\mathrm{Numerical}}-Q_{\mathrm{Analytical}}}{Q_{\mathrm{Analytical}}}
\right|
\label{eq:mesh_relative_error}
\end{equation}
\end{linenomath*}

As illustrated in Figure~S5, the relative flow-rate errors for all mesh scales
are strictly confined within 1.5\%. The error exhibits a clear downward trend
as the mesh resolution increases ($\Delta x$ decreases): the relative errors for
$N_Y=100$, $200$, $400$, and $800$ are 1.2995\%, 1.2165\%, 0.8230\%, and
0.7380\%, respectively. Considering the balance between computational accuracy
and cost, the range $N_Y=200$--$400$ is selected as the production grid for
this study.


%% -----------------------------------------------------------------------
%%  TEXT S3 -- Trapping Characterization and Trapping Criterion Validation
%% -----------------------------------------------------------------------
\noindent\textbf{Text S3.}
\textbf{Validation of the trapping criterion}

Detailed extraction and connectivity procedures are provided in the main text.
Here we retain only the trapping-classification criterion used in the validation.
A cluster was classified as trapped when its transport activity fell below a
capillary threshold:
\begin{linenomath*}
\begin{equation}
\label{eq:locking_criterion}
\overline{U}_x^{(k)}
\le
C_a^{*}\,
\frac{\sigma}{\mu_{\phi}}
\end{equation}
\end{linenomath*}
where $C_a^{*}$ is a critical capillary number, $\sigma$ denotes the gas--water
interfacial tension, and $\mu_{\phi}$ is the viscosity of the corresponding fluid
phase. This criterion encapsulates the regime where viscous driving forces are
insufficient to overcome interfacial energy barriers, resulting in negligible net
transport within the cluster.

We performed a two-parameter scan around the calibrated trapping criterion
(phase A = water, phase B = gas),
including the base threshold $C_{a,A}^{*}$ and the phase-B scaling factor
$\lambda_B$ (i.e., $C_{a,B}^{*}=\lambda_B C_{a,A}^{*}$). For each parameter pair,
the phase trapping indices from all simulations were recalculated and compared
with the corresponding relative permeabilities.

Consistency was evaluated primarily using Spearman correlation. Let $\eta_A$
and $\eta_B$ denote the phase trapping indices, and
$k_{rw}$ and $k_{rg}$ denote water and gas relative permeabilities, respectively.
The Spearman function is:
\begin{linenomath*}
\begin{equation}
\label{eq:spearman_def}
\rho_s(X,Y)
=
\rho_p\!\left(\mathrm{rank}(X),\,\mathrm{rank}(Y)\right)
=
1-\frac{6\sum_{i=1}^{n} d_i^2}{n\left(n^2-1\right)}
\end{equation}
\end{linenomath*}
where $d_i$ is the rank difference between paired samples (for tie-free data).
Physically, $\rho_s$ measures monotonic coupling between trapping state and
hydraulic conductance.

The scan diagnostics are:
\begin{linenomath*}
\begin{equation}
\label{eq:scan_corr_A}
S_A = \left|\rho_s\!\left(\eta_A,\,k_{rw}\right)\right|,
\qquad
S_B = \left|\rho_s\!\left(\eta_B,\,k_{rg}\right)\right|
\end{equation}
\end{linenomath*}
where $\rho_s(\cdot)$ is the Spearman rank correlation coefficient. Pearson
coefficients were reported in parallel only as secondary linear-consistency
checks and were not used as the scan objective.
To avoid selecting a parameter set that performs well for only one phase,
we used a balanced score:
\begin{linenomath*}
\begin{equation}
\label{eq:scan_hmean_score}
S_{\mathrm{bal}}
=
\frac{2S_A S_B}{S_A+S_B}
\end{equation}
\end{linenomath*}
which penalizes one-sided performance.

Figure~S6 summarizes the parameter-scan map. High-score regions are concentrated
in a narrow neighborhood of the adopted parameter set, and the optimum is stable
across adjacent grid points. The unscaled criterion shows an inter-phase mismatch:
the parameter region maximizing phase-A consistency does not fully coincide with
that of phase B. Introducing the phase-B scaling parameter $\lambda_B$ reduces
this mismatch and yields a robust two-phase optimum, suggesting that the effective
trapping threshold is not strictly symmetric between gas and water.

To further verify the selected optimum, Figure~S7 reports direct pointwise
correlations between trapping index and relative permeability at
$C_a^* = 1.8 \times 10^{-3}$ and $\lambda_B = 2.333$ for $n=130$ samples
(pore-filling morphology, and grain-coating morphology,combined). Both phases exhibit the expected negative
direction ($\eta \uparrow \Rightarrow k_r \downarrow$): for phase A,
$\rho_s=-0.960$ and $r=-0.741$; for phase B, $\rho_s=-0.820$ and $r=-0.527$.
The weaker linear correlation in phase B, compared with its still-strong rank
correlation, is consistent with a threshold-like mobility collapse. Given the
strongly nonlinear and jump-like mobility response of the gas phase, Spearman
rank correlation provides a more objective measure of the physical association
between $\eta$ and $k_r$ than Pearson correlation. Taken together,
Figures~S6 and~S7 demonstrate that $\eta$ provides a robust physical proxy for
phase mobility loss in hydrate-bearing sediments.
The trapping criterion calibrated on the 130-case subset
was applied uniformly to the full 220-case
SHAP ensemble; consistency was verified.

%% -----------------------------------------------------------------------
%%  TEXT S4 -- Simulation Design and SHAP Analysis Strategy
%% -----------------------------------------------------------------------
\noindent\textbf{Text S4.}
\textbf{Model Validation and SHAP Attribution Analysis}

XGBoost models are trained to predict phase trapping indicators.
After retaining $n=180$ samples with $S_w\in(0,1)$, each target variable
is evaluated using a 90\%/10\% train/holdout split, where the holdout
subset is set aside before any model fitting or selection.

Two strategies are compared: a baseline XGBoost with fixed hyperparameters
and a tuned XGBoost whose hyperparameters are optimized on the training
subset via RandomizedSearchCV\@. For each strategy, 5-fold cross-validation
is performed on the training subset, and the final model is retrained on
all training data and assessed on the holdout subset.
Metrics are the coefficient of determination $R^2$ and mean absolute
error (MAE).

Cross-validation fold results are reported in Figure~S8, and holdout
parity plots in Figure~S9. For the water-phase indicator, the baseline
XGBoost is adopted as the final model (holdout $R^2=0.914$,
MAE$\approx 0.085$; cross-validation mean $\overline{R^2}=0.6927\pm0.1066$).
For the gas-phase indicator, the tuned XGBoost is selected (holdout
$R^2=0.755$, MAE$\approx 0.152$; cross-validation mean
$\overline{R^2}=0.5963\pm0.0623$).

This protocol reduces potential data leakage and model-selection bias,
though the limited sample size and moderate cross-validation variability
indicate that the predictive uncertainty is non-negligible.
The validated models are used in the main-text SHAP attribution analysis,
where feature contributions to the water- and gas-phase trapping indicators
are quantified to identify the dominant controlling factors.
Future work could strengthen the analysis by incorporating larger
simulation ensembles, extending feature sets with additional geometric
descriptors, and adopting model-agnostic importance measures to
complement the SHAP attribution framework.

%% -----------------------------------------------------------------------
%%  REFERENCES
%% -----------------------------------------------------------------------
% NOTE: All references listed here must also appear in the main manuscript.
% Use only \cite and \citeA for citations in AGU SI documents.
% \cite for parenthetical: ...as shown in recent studies \cite{lallemand_theory_2000}
% \citeA for in-text:       \citeA{lallemand_theory_2000} have shown...

\bibliography{reference_si}

\end{article}
\clearpage

%% -----------------------------------------------------------------------
%%  FIGURES  (must appear after \end{article})
%% -----------------------------------------------------------------------

% Figure S1 -- Laplace pressure test (linear DeltaP vs 1/R)

\begin{figure}
\setfigurenum{S1}
\centering
\includegraphics[width=0.6\textwidth]{fig1_laplace_eq1.pdf}\\[-0.2em]
{\ (a) Eq1}\\[0.8em]
\includegraphics[width=0.6\textwidth]{fig1_laplace_eq2.pdf}\\[-0.2em]
{\ (b) Eq2}
\caption{Laplace pressure test. Linear fits
($R^2 \approx 1$) confirm that the model reproduces the Laplace relation
$\Delta P = \sigma/R$ and yields reliable estimates of interfacial tension.
Eq1 and Eq2 denote the two thermodynamic equilibrium states listed in Table~S1.}\label{fig:S1_laplace_eq1}
\end{figure}

% Figure S2 -- Contact angle vs adhesive coefficient
\begin{figure}
\setfigurenum{S2}
\centering
\includegraphics[width=0.6\textwidth]{fig2_contact_angle_eq1.pdf}\\[-0.2em]
{\ (a) Eq1}\\[0.8em]
\includegraphics[width=0.6\textwidth]{fig2_contact_angle_eq2.pdf}\\[-0.2em]
{\ (b) Eq2}
\caption{Contact angle validation. Equilibrium droplet configurations under both
equilibrium states (Eq1 and Eq2) for varying adhesive coefficients $G_{k,\mathrm{ads}}$.
A linear relationship between $G_{k,\mathrm{ads}}$ and the resulting static contact
angle $\theta$ is observed, confirming that the adhesive coefficient directly controls
surface wettability. Eq1 and Eq2 denote the two thermodynamic equilibrium states
listed in Table~S1.}\label{fig:S2_contact_angle}
\end{figure}





% Figure S3 -- Schematic of two-dimensional two-phase Poiseuille flow
\begin{figure}
\setfigurenum{S3}
\centering
\includegraphics[width=0.92\textwidth,trim=5cm 6cm 5cm 6cm,clip]{fig3_two_phase_flow.pdf}
\caption{Schematic of the two-dimensional two-phase Poiseuille flow.
Wetting and non-wetting regions, interface positions ($\pm a$), and channel
half-width ($\pm b$) are indicated together with the flow direction and
coordinate system used in the analytical derivation.}\label{fig:S3_poiseuille_schematic}
\end{figure}

% Figure S4 -- Two-phase Poiseuille flow relative-permeability validation
\begin{figure}
\setfigurenum{S4}
\centering
\includegraphics[width=0.6\textwidth]{fig4_relperm_eq1.pdf}\\[-0.2em]
{\ (a) Eq1}\\[0.8em]
\includegraphics[width=0.6\textwidth]{fig4_relperm_eq2.pdf}\\[-0.2em]
{\ (b) Eq2}
\caption{Two-dimensional two-phase Poiseuille flow validation. Simulated
relative permeabilities for wetting and non-wetting phases agree with the
analytical solutions in both equilibrium states ($R^2 \approx 1$). Eq1 and Eq2
denote the two thermodynamic equilibrium states listed in Table~S1.}\label{fig:S4_relperm}
\end{figure}

% Figure S5 -- Mesh independence validation
\begin{figure}
\setfigurenum{S5}
\centering
\includegraphics[width=0.72\textwidth]{fig5_mesh_convergence.pdf}
\caption{Mesh-independence validation based on  Poiseuille-flow
benchmarks. Relative flow-rate error decreases with increasing resolution
($N_Y=100, 200, 400, 800$), and remains below 1.5\% for all tested grids,
supporting the choice of $N_Y=200$--$400$ for production simulations.}\label{fig:S5_mesh_convergence}
\end{figure}


\begin{figure}
\setfigurenum{S6}
\centering
\includegraphics[width=0.95\textwidth]{fig6_trapping_validation.pdf}
\caption{Parameter-scan validation of trapping-index consistency with relative permeability.
Each panel reports the absolute Spearman correlation for one phase and the
harmonic-mean balanced score, as functions of $C_{a,A}^{*}$ and $\lambda_B$.
The concentration of high-score regions around the selected parameter set
demonstrates robust two-phase consistency and supports the adopted calibration.}\label{fig:S6_trapping_validation}
\end{figure}

\begin{figure}
\setfigurenum{S7}
\centering
\includegraphics[width=0.95\textwidth]{fig7_eta_vs_kr_scatter.pdf}
\caption{{Correlation between phase-trapping metric $\eta$ and relative permeability $k_r$
for (a) water phase and (b) gas phase.} Dots and squares represent Morph-1 and Morph-2
results, respectively. The black dashed line indicates the linear regression with a
95\% confidence interval (shaded area). Statistical metrics including Pearson $r$,
Spearman $\rho$, and sample size $n$ are provided in the legends.}\label{fig:S7_trapping_validation}
\end{figure}



\begin{figure}
\setfigurenum{S8}
\centering
\includegraphics[width=0.6\textwidth]{fig8_K_validation_gas.pdf}\\[-0.2em]
{\ (a)  Cross-Validation for Gas Phase}\\[0.8em]
\includegraphics[width=0.6\textwidth]{fig8_K_validation_water.pdf}\\[-0.2em]
{\ (b) Cross-Validation for Water Phase}
\caption{Five-fold outer-cross-validation results for XGBoost models
predicting phase trapping indicators.}\label{fig:S8_cross_validation}
\end{figure}


\begin{figure}
\setfigurenum{S9}
\centering
\includegraphics[width=0.49\textwidth]{fig9_validation_gas_parity.pdf}
\hfill
\includegraphics[width=0.49\textwidth]{fig9_validation_water_parity.pdf}\\[-0.2em]
{\ (a) Gas-phase indicator (tuned XGBoost)
\hspace{1.5em}
(b) Water-phase indicator (baseline XGBoost)}
\caption{Holdout parity plots of selected XGBoost models for phase trapping indicators.
Markers show predicted versus observed values on the independent holdout subset,
and the dashed diagonal denotes the 1:1 agreement line.}\label{fig:S9_holdout_parity}
\end{figure}



\begin{figure}
\setfigurenum{S10}
\centering
\includegraphics[width=0.5\textwidth]{fig10_eq_effect_bars.pdf}\\[-0.2em]
{\ (a) Pairwise SHAP comparison by equilibrium state (Eq1 vs Eq2)}\\[0.2em]
\includegraphics[width=0.5\textwidth]{fig10_morph_effect_bars.pdf}\\[-0.2em]
{\ (b) Pairwise SHAP comparison by pore morphology (Morph-1 vs Morph-2)}
\includegraphics[width=0.5\textwidth]{fig10_shap_gx_grouped_bars.pdf}\\[-0.2em]
{\ (c) Pairwise SHAP comparison by forcing factor ($G_x$ groups)}
\caption{Supplementary pairwise SHAP comparisons for phase trapping indicators.
Panels (a)--(c) compare phase-wise SHAP responses across equilibrium state,
pore morphology, and forcing factor groupings, respectively. These plots
complement the main-text phase-wise SHAP response curves by showing additional
grouped contrasts beyond the four highlighted indicators.}\label{fig:S10_shap_pairwise}
\end{figure}




%% -----------------------------------------------------------------------
%%  TABLES  (may appear after \end{article})
%% -----------------------------------------------------------------------

% Table S1 -- Equilibrium states
\begin{table}
\settablenum{S1}
\caption{Equilibrium states with corresponding pressure--temperature conditions
and derived fluid properties used in the simulations.}
\centering
\renewcommand{\arraystretch}{1.6}
\setlength{\tabcolsep}{9pt}
\begin{tabular}{lccccccc}
\toprule
State &
$P$ &
$T$ &
$\rho_w$ &
$\mu_w$ &
$\rho_g$ &
$\mu_g$ &
$\sigma$ \\
&
(MPa) &
(K) &
(kg\,m$^{-3}$) &
(mPa\,s) &
(kg\,m$^{-3}$) &
(mPa\,s) &
(mN\,m$^{-1}$) \\
\midrule
Eq1 & 7.270 & 283.1 & 1003 & 1.302 & 57.82 & 0.01237 & 63.91 \\
Eq2 & 4.300 & 278.2 & 1002 & 1.510 & 32.91 & 0.01126 & 66.68 \\
\bottomrule
\multicolumn{8}{p{0.95\linewidth}}{\footnotesize
\textit{Note.} $P$ and $T$ denote pressure and temperature, respectively;
$\rho_w$ and $\mu_w$ are the density and dynamic viscosity of water;
$\rho_g$ and $\mu_g$ are the density and dynamic viscosity of methane;
$\sigma$ is the gas--water surface tension.}
\end{tabular}
\end{table}

% Table S2 -- Scaling parameters
\begin{table}
\settablenum{S2}
\caption{Simulation scaling parameters and lattice-to-physical unit conversion relations.}
\centering
\renewcommand{\arraystretch}{2.5}
\setlength{\tabcolsep}{9pt}
\begin{tabular}{P{6.0cm}P{3.0cm}P{6cm}}
\toprule
Quantity & L.U. & P.U. / Reference expression \\
\midrule
Lattice spacing $\delta_x$ & 1.0 & $X_r = 2.5 \times 10^{-7}$ m \\
Time step $\delta_t$       & 1.0 & $T_r = 1.5 \times 10^{-9}$ s \\
Reference density $\rho_r$ & --- &
  \makecell[l]{Eq1: 132.86 \\ Eq2: 138.70} \\
Water density $\rho_w$ &
  \makecell[l]{7.5502 (Eq1)\\ 7.2244 (Eq2)} &
  scaled as $\rho_r$ \\
Methane density $\rho_g$ &
  \makecell[l]{0.4352 (Eq1)\\ 0.2373 (Eq2)} &
  scaled as $\rho_r$ \\
Water kinematic viscosity $\nu_w$ &
  \makecell[l]{0.03113 (Eq1)\\ 0.03613 (Eq2)} &
  scaled as $X_r^{2} T_r^{-1}$ \\
Methane kinematic viscosity $\nu_g$ &
  \makecell[l]{0.005133 (Eq1)\\ 0.0082 (Eq2)} &
  scaled as $X_r^{2} T_r^{-1}$ \\
Surface tension $\sigma$ &
  \makecell[l]{0.06927 (Eq1)\\ 0.06923 (Eq2)} &
  scaled as $\rho_r X_r^{3} T_r^{-2}$ \\
\bottomrule
\end{tabular}
\end{table}


\end{document}
